% problem_id: S001-lemma-envelope
% source_id: S001 (Stacks Project)
% source_file: chow.tex
% statement_locator: chow.tex lines 3592--3605
% proof_locator: chow.tex lines 3607--3709
% env: lemma
% id_note: label 在本来源内唯一
% pairing: tight (verified)
% status: complete (statement + author proof verbatim + reason + locators)
%
\begin{lemma}
\label{lemma-envelope}
Let $(S, \delta)$ be as in Situation \ref{situation-setup}.
Let $X$ be a scheme locally of finite type over $S$.
Let $f : Y \to X$ be an envelope. Then
we have an exact sequence
$$
\CH_k(Y \times_X Y) \xrightarrow{p_* - q_*}
\CH_k(Y) \xrightarrow{f_*}
\CH_k(X) \to 0
$$
for all $k \in \mathbf{Z}$. Here $p, q : Y \times_X Y \to Y$ are
the projections.
\end{lemma}

\begin{proof}
Since $f$ is an envelope, $f$ is proper and hence pushforward on
cycles and cycle classes is defined, see
Sections \ref{section-proper-pushforward} and \ref{section-push-pull}.
Similarly, the morphisms $p$ and $q$ are proper as base changes of $f$.
The composition of the arrows is zero as
$f_* \circ p_* = (p \circ f)_* = (q \circ f)_* = f_* \circ q_*$, see
Lemma \ref{lemma-compose-pushforward}.

\medskip\noindent
Let us show that $f_* : Z_k(Y) \to Z_k(X)$ is surjective.
Namely, suppose that we have $\alpha = \sum n_i[Z_i] \in Z_k(X)$
where $Z_i \subset X$ is a locally finite family of integral
closed subschemes. Let $x_i \in Z_i$ be the generic point.
Since $f$ is an envelope and hence completely decomposed,
there exists a point $y_i \in Y$ with $f(y_i) = x_i$
and with $\kappa(y_i)/\kappa(x_i)$ trivial. Let $W_i \subset Y$
be the integral closed subscheme with generic point $y_i$.
Since $f$ is closed, we see that $f(W_i) = Z_i$.
It follows that the family of closed subschemes $W_i$ is locally finite
on $Y$. Since $\kappa(y_i)/\kappa(x_i)$ is trivial we see
that $\dim_\delta(W_i) = \dim_\delta(Z_i) = k$. Hence
$\beta = \sum n_i[W_i]$ is in $Z_k(Y)$. Finally, since
$\kappa(y_i)/\kappa(x_i)$ is trivial, the degree of the dominant
morphism $f|_{W_i} : W_i \to Z_i$ is $1$ and we conclude
that $f_*\beta = \alpha$.

\medskip\noindent
Since $f_* : Z_k(Y) \to Z_k(X)$ is surjective, a fortiori the map
$f_* : \CH_k(Y) \to \CH_k(X)$ is surjective.

\medskip\noindent
Let $\beta \in Z_k(Y)$ be an element such that $f_*\beta$ is zero in
$\CH_k(X)$. This means we can find a locally finite family of
integral closed subschemes $Z_j \subset X$ with $\dim_\delta(Z_j) = k + 1$
and $f_j \in R(Z_j)^*$ such that
$$
f_*\beta = \sum (Z_j \to X)_*\text{div}(f_j)
$$
as cycles where $i_j : Z_j \to X$ is the given closed immersion.
Arguing exactly as above, we can find a locally finite
family of integral closed subschemes $W_j \subset Y$
with $f(W_j) = Z_j$ and such that $W_j \to Z_j$ is birational, i.e.,
induces an isomorphism $R(Z_j) = R(W_j)$. Denote $g_j \in R(W_j)^*$
the element corresponding to $f_j$. Observe that $W_j \to Z_j$
is proper and that $(W_j \to Z_j)_*\text{div}(g_j) = \text{div}(f_j)$
as cycles on $Z_j$. It follows from this that if we replace
$\beta$ by the rationally equivalent cycle
$$
\beta' = \beta - \sum (W_j \to Y)_*\text{div}(g_j)
$$
then we find that $f_*\beta' = 0$.
(This uses Lemma \ref{lemma-compose-pushforward}.)
Thus to finish the proof
of the lemma it suffices to show the claim in the following paragraph.

\medskip\noindent
Claim: if $\beta \in Z_k(Y)$ and $f_*\beta = 0$, then
$\beta = \delta + p_*\gamma - q_*\gamma$ in $Z_k(Y)$ for some
$\gamma \in Z_k(Y \times_X Y)$. Namely, write $\beta = \sum_{j \in J} n_j[W_j]$
with $\{W_j\}_{j \in J}$ a locally finite family of integral closed
subschemes of $Y$ with $\dim_\delta(W_j) = k$.
Fix an integral closed subscheme $Z \subset X$. Consider the subset
$J_Z = \{j \in J : f(W_j) = Z\}$. This is a finite set. There are three
cases:
\begin{enumerate}
\item $J_Z = \emptyset$. In this case we set $\gamma_Z = 0$.
\item $J_Z \not = \emptyset$ and $\dim_\delta(Z) = k$.
The condition $f_*\beta = 0$ implies by looking at the
coefficient of $Z$ that $\sum_{j \in J_Z} n_j\deg(W_j/Z) = 0$.
In this case we choose an integral closed subscheme $W \subset Y$
which maps birationally onto $Z$ (see above). Looking at generic
points, we see that $W_j \times_Z W$ has a unique irreducible
component $W'_j \subset W_j \times_Z W \subset Y \times_X Y$
mapping birationally to $W_j$. Then $W'_j \to W$ is dominant
and $\deg(W'_j/W) = \deg(W_j/W)$. Thus if we set
$\gamma_Z = \sum_{j \in J_Z} n_j[W'_j]$
then we see that
$p_*\gamma_Z = \sum_{j \in J_Z} n_j[W_j]$ and
$q_*\gamma_Z = \sum_{j \in J_Z} n_j\deg(W'_j/W)[W] = 0$.
\item $J_Z \not = \emptyset$ and $\dim_\delta(Z) < k$.
In this case we choose an integral closed subscheme $W \subset Y$
which maps birationally onto $Z$ (see above). Looking at generic
points, we see that $W_j \times_Z W$ has a unique irreducible
component $W'_j \subset W_j \times_Z W \subset Y \times_X Y$
mapping birationally to $W_j$. Then $W'_j \to W$ is dominant
and $k = \dim_\delta(W'_j) > \dim_\delta(W) = \dim_\delta(Z)$.
Thus if we set $\gamma_Z = \sum_{j \in J_Z} n_j[W'_j]$
then we see that
$p_*\gamma_Z = \sum_{j \in J_Z} n_j[W_j]$ and
$q_*\gamma_Z = 0$.
\end{enumerate}
Since the family of integral closed subschemes $\{f(W_j)\}$
is locally finite on $X$
(Lemma \ref{lemma-quasi-compact-locally-finite})
we see that the $k$-cycle
$$
\gamma = \sum\nolimits_{Z \subset X\text{ integral closed}} \gamma_Z
$$
on $Y \times_X Y$ is well defined. By our computations above it follows that
$p_*\gamma_Z = \beta$ and $q_*\gamma_Z = 0$ which implies
what we wanted to prove.
\end{proof}
